% Discussion paragraphs from the mnist_holes / 33-language thread; additive to discussion.tex.
\subsection{What reasoning buys, and where it runs out}
\paragraph{What reasoning buys.} The text-quality and hole-finder networks were written in a few hours of
conversation and have no training cost, and both land at the level of the classical
hand-engineered baselines for their tasks: the text gate reproduces most of a production
filtering pipeline's decisions across 33 languages, and the second digit classifier beats a trained prototype classifier. Both are entirely legible: every score decomposes into named
contributions, which is what a data engineer wants from a filter (``why was this document
dropped?'') and what makes the error analyses above possible without any probing.

\paragraph{Where reasoning runs out.} The failures are informative. The text network
cannot see grammar, because that needs a representation of words and their order, which
no amount of character-level reasoning produces; and the 33-language ablation shows that
several detectors that are easy to justify contribute nothing measurable. The digit
network's ceiling is set by the resolution of what a person can specify: a $5\times5$
sketch cannot say ``the loop must close'', and the hole finder cannot close a loop the
writer left open. A trained network learns the fine stroke-topology statistics that fill
that gap, and that is exactly the part of the task a designer cannot write down.

\paragraph{Fine-tuning from a designed initialisation.} Because these are ordinary networks
with PyTorch exports (\texttt{to\_torch()}, checked against NumPy in the tests), the
obvious next experiment is to use them as initialisations: does a network that starts
from named, correct detectors train faster, need fewer labels, or stay legible? The
repository deliberately does not run that experiment yet, so that the reported numbers
are those of the hand-written weights alone.

\paragraph{On process.} The weights were written by a language model in conversation with
the first author, who set the tasks, the rule (reason and check, never fit) and the
standard of evidence. We think this is worth stating because it bears on what the result
means: the design knowledge in these networks is the kind that can be articulated---what
mojibake looks like, where a 4 is empty---and a system that can articulate it can also
write it into weights. The repository documents the process, including the development
log of Section~\ref{sec:tq-devlog} and the iteration on training-set confusion matrices
for MNIST, so that the line between ``reasoned'' and ``fit'' is visible to the reader.

